% Proved post-CP22 example. This is an obstruction to automatic tameness only.
\subsection{A wildly ramified cyclic cover with one smooth exceptional curve}

\begin{example}[Proved: failure of automatic radial tameness]
\label{D:one-curve-wild-swap-cover}
Let $k$ be any algebraically closed field of characteristic $2$, and
let $C/k$ be a smooth projective connected curve of genus $g\geq2$.
There is a complete excellent normal local $k$-domain $R$ of dimension
two, with residue field $k$, and a projective regular resolution
\[
       \pi_R:Y_R\longrightarrow\operatorname{Spec}R
                                                        \tag{D.WS.1}
\]
whose reduced exceptional divisor is a single smooth curve $D$ of
genus $g$ and square $D^2=1-g$. Its puncture
$U=\operatorname{Spec}R\setminus\{\mathfrak m_R\}$ has a connected
finite \'etale $\mathbb Z/2$-torsor which is not radially tame along
$D$. In the corresponding extension of divisorial discrete valuation
rings, the ramification index is $1$ and the residue extension is
purely inseparable of degree $2$. The cover does not extend to a
finite \'etale cover of $Y_R$.

Consequently the hypothesis of a radially tame outer or permutation
quotient in the one-curve crown results does not follow merely from
the existence of one smooth exceptional curve. This example makes
no claim against the representation-count bounds or finite generation.
\end{example}

\begin{proof}
Put $W=C\times C$, let $\sigma$ interchange the two factors, and let
\[
       f:W\longrightarrow Y=\operatorname{Sym}^2 C
                                                        \tag{D.WS.2}
\]
be the degree-two quotient. Write $\Delta\subset W$ for the diagonal
and $D\subset Y$ for its reduced image. We include the local
calculation because ramification index alone does not detect the
wild ramification here.

At a point $(c,c)$, choose the same uniformizer on both factors.
The completed local quotient is
\[
 \begin{aligned}
 k[[x,y]]^{\langle\sigma\rangle}&=k[[u,v]],
       &u=x+y,\quad v=xy,\\
 k[[x,y]]&=k[[u,v]][x]/(x^2+ux+v).
 \end{aligned}                                           \tag{D.WS.3}
\]
The symmetric-power-series identity follows from the elementary
symmetric-polynomial identity by completion; the two adic topologies
are cofinal by the displayed monic relation. Finite-group invariants
commute with this flat completion, since they are the kernel of the
finite collection of maps $\sigma-1$. The right-hand presentation is
finite free with basis $1,x$ and identifies the complete local rings.
Away from $\Delta$ the constant group $\mathbb Z/2$ acts freely, so
the quotient is \'etale there. These descriptions show that $Y$ is a
smooth projective surface and that $f$ is finite flat of degree two.

The equation of $D$ is $u=0$, whereas that of $\Delta$ is $x+y=0$.
In particular
\[
       f^*D=\Delta,\qquad
       f_\Delta:\Delta\longrightarrow D
       \text{ is relative Frobenius of degree }2.
                                                        \tag{D.WS.4}
\]
Indeed modulo $u$ the map is $k[[v]]\to k[[x]]$, $v\mapsto x^2$.
The same description of the generic residue fields identifies $D$
with the Frobenius twist of $C$; thus $D$ is smooth of genus $g$.

We construct the contracting line bundle on $Y$ itself. Choose an
effective canonical divisor $K_C=\sum_c n_c c$, and let
$X_c=\{c+z:z\in C\}\subset\operatorname{Sym}^2C$. Define
\[
                  H=\sum_c n_c X_c,\qquad L=H+2D.
                                                        \tag{D.WS.5}
\]
Each $X_c$ is Cartier on the smooth surface, and divisor pullback gives
\[
            f^*H=\operatorname{pr}_1^*K_C+
                         \operatorname{pr}_2^*K_C.
                                                        \tag{D.WS.6}
\]
This also holds at $2c$: there $X_c$ has equation $v=0$, whose
pullback is $xy=0$. Thus this construction uses an actual divisor;
it does not assume descent of a line bundle through a wild quotient
from triviality of stabilizer actions. The pullback of $H$ is ample,
because $\deg K_C=2g-2>0$, and finite surjective descent of ampleness
makes $H$ ample.

The normal bundle of $\Delta$ is $T_C$. Intersection theory and
(D.WS.4)--(D.WS.6) therefore give
\[
 \begin{aligned}
 \Delta^2&=2-2g,& D^2&=1-g,\\
 H\cdot D&=2g-2,& H^2&=4(g-1)^2,\\
 L\cdot D&=0,& L^2&=4g(g-1)>0.
 \end{aligned}                                           \tag{D.WS.7}
\]
For every irreducible curve $B\ne D$, the two terms
$H\cdot B>0$ and $D\cdot B\geq0$ imply $L\cdot B>0$.
Hence $L$ is nef and big, with null locus exactly $D$.
Moreover
\[
 f_\Delta^*(L|_D)
       \simeq\omega_C^{\otimes2}\otimes T_C^{\otimes2}
       \simeq\mathcal O_C.
                                                        \tag{D.WS.8}
\]
The finite map $f_\Delta$ is flat of degree two between smooth
curves. Applying its norm to (D.WS.8) gives
$(L|_D)^{\otimes2}\simeq\mathcal O_D$. Thus $L|_D$ is semiample.

\citet[Theorem0.2, or directly Proposition1.6 with the effective
Cartier divisor $2D$]{C_Keel1999WildSwap} now implies that $L$ is
semiample. The direct application uses $L=H+2D$ with $H$ ample;
Proposition1.6 and its descent lemma are proved on pp.~261--264.
The related contraction of the diagonal of $C\times C$ by
$\omega_{C\times C}(2\Delta)$ is already discussed in
\citet[Section3, Theorem3.0 and Lemma3.2]{C_Keel1999WildSwap}.
We use that construction here only to produce the stated wild-cover
obstruction, and make no novelty claim for the diagonal contraction.

Take a basepoint-free positive power of $L$, its morphism to the
projective image, and its normal Stein factor. This gives a
projective birational morphism
\[
                  \pi:Y\longrightarrow X
                                                        \tag{D.WS.9}
\]
to a normal projective surface, with $\pi_*\mathcal O_Y=\mathcal O_X$.
Its contracted curves are exactly the curves of $L$-degree zero,
so $D$ is the only one. Birationality follows equally from bigness
and the normal Stein factor. Proper birational quasi-finite maps
to a normal scheme are isomorphisms; thus $\pi$ is an isomorphism
off $D$. Its connected fibres show that the set-theoretic fibre
over $x=\pi(D)$ is exactly $D$.

Apply Stein factorization once more, now to $\pi f$:
\[
             W\xrightarrow{\rho} Z\xrightarrow{h}X.
                                                        \tag{D.WS.10}
\]
Here $Z$ is normal, $h$ is finite of generic degree two, and
$\rho$ has geometrically connected fibres; see
\citet[Tag03H0, complete Noetherian Stein-factorization theorem]
{C_StacksWildSwap}. Over $X\setminus\{x\}$ this is the original
finite map $f$, so $\rho$ is an isomorphism there and $h$ is a
$\mathbb Z/2$-torsor there. The set-theoretic fibre of $\pi f$
over $x$ is $\Delta$. It is connected and surjects onto the finite
set $h^{-1}(x)$, so that set consists of a single point $z$.

Set
\[
       R=\widehat{\mathcal O}_{X,x},\qquad
       A=\widehat{\mathcal O}_{Z,z}.
                                                        \tag{D.WS.11}
\]
The local rings before completion are excellent normal local
domains essentially of finite type over $k$. Their completions
are normal local domains; in particular they do not split into
several branches. These are the normal-completion consequences
of excellence in \citet[Tags07QS,0C23 and0C22]
{C_StacksWildSwap}. Since $z$ is the unique point above $x$,
completion of the finite algebra gives exactly the local
$R$-algebra $A$, rather than a product of several completions.
The degree remains two. The involution extends through Stein
factorization and completion, and its invariants are $R$:
flat completion commutes with the kernel defining invariants.
The finite map
\[
 \operatorname{Spec}A\setminus\{\mathfrak m_A\}
                 \longrightarrow U
                                                        \tag{D.WS.12}
\]
is therefore \'etale of degree two with the given involution.
Its source is a nonempty open of the integral scheme
$\operatorname{Spec}A$, hence is connected. Thus (D.WS.12) is
a connected $\mathbb Z/2$-torsor, not a cover which splits after
completion.

Base change of (D.WS.9) to $\operatorname{Spec}R$ gives (D.WS.1).
The completion map is regular by excellence. Its finite-type
base changes are regular, and a regular map over a regular
scheme has regular source; see \citet[Tags07C1 and07QI]
{C_StacksWildSwap}. Hence $Y_R$ is regular. Its exceptional
fibre is still the single smooth curve $D$, with the same
normal bundle and self-intersection. Similarly $W_R$ is regular.
The finite map $W_R\to Y_R$ is the normalization in the function
field of (D.WS.12). Its divisor above $D$ is still $\Delta$,
and (D.WS.4) is unchanged by this base change. At their generic
points it follows that
\[
      e(\Delta/D)=1,\qquad
      [k(\Delta):k(D)]_{\rm insep}=2.
                                                        \tag{D.WS.13}
\]
The full involution acts trivially on the residue field, so
the inertia group has order two. This is fierce wild
ramification: tameness also requires a separable residue
extension, not just a prime-to-characteristic ramification
index. It excludes extension of the punctured torsor to a
finite \'etale cover of $Y_R$, since the finite normal
extension is unique.

The Artin--Schreier description can be seen without an
inertia convention. In (D.WS.3), after inverting $u$, put
$w=x/u$. Then
\[
                        w^2+w=\frac{v}{u^2}.
                                                        \tag{D.WS.14}
\]
In the completed divisorial field the leading residue $v$
is not a square in $k((v))$. An Artin--Schreier correction
of pole order one could cancel the pole of order two only
if its leading coefficient squared to $v$; larger pole
orders produce a strictly larger uncancellable leading
pole. Thus this local equation also displays the failure
of an unramified extension.

Finally, this is genuinely an obstruction already among
$p$-cyclic classes. Normality and proper birationality give
$\Gamma(U,\mathcal O_U)=\Gamma(Y_R,\mathcal O_{Y_R})=R$.
The map $a\mapsto a^2-a$ is surjective on the strictly
henselian ring $R$. The Artin--Schreier exact sequence
therefore identifies $H^1(-,\mathbb F_2)$ on both spaces
with the Frobenius-fixed part of $H^1(-,\mathcal O)$.
The class of (D.WS.12) is outside the image of
$H^1(Y_R,\mathbb F_2)\to H^1(U,\mathbb F_2)$, by
(D.WS.13). Thus finiteness of the latter cohomology group
does not force its classes to extend across $D$.
\end{proof}
